\section{Introduction}\label{sec:intro}

Multi-object tracking (MOT) finds every object of interest in a video and
gives each object the same identity over time. It is used in surveillance,
driving, and drone systems
\cite{zhu2022visdrone,wu2022uavsurvey}, and many of these
systems must give results with little delay. Most trackers follow the
tracking-by-detection design \cite{adzemovic2026survey}. A detector gives a score to each candidate box
in a frame, and the tracker links the boxes whose scores pass its thresholds
into tracks \cite{bewley2016sort,zhang2022bytetrack,cao2023ocsort}.

In real systems, the detector and the tracker are rarely built together. A
detector is replaced by a faster model, retrained for a new camera, or made
smaller (quantized) for embedded hardware. The tracker, however, keeps the
thresholds that its authors set for another detector. Detector scores are not
calibrated across models and conditions \cite{kuppers2020detcalib}, so a new
detector usually gives scores on a different scale, and nothing in the system
notices the change.

The effect can be large. In this paper, halving the scores of a published
detector reduces two well-known trackers to 0 higher order tracking accuracy
(HOTA) \cite{luiten2021hota}, because no detection passes their fixed
thresholds (Sect.~\ref{sec:shift}). Retuning every tracker for every detector
needs labelled data from the new setting, which a working system usually does
not have.

Current methods solve only part of this problem. Trackers with adaptive
thresholds change their own thresholds from the detections of each frame, but
each method is built into one tracker
\cite{ma2023adaptivebytetrack,shim2024adaptrack,act2026adaptive}. Score
calibration methods correct the detector scores, but they need labelled data
\cite{kuppers2020detcalib}. Adaptive inference methods learn or measure a
policy for one model \cite{han2022dynamic,jiang2018chameleon}. What is missing
is a component that needs no training, works outside the tracker, adjusts itself
from the detector's own scores, and leaves the tracker alone when it already
fits those scores.

This paper presents such a component: a self-calibrating control layer (\scl{})
that sits between a fixed (frozen) detector and a fixed tracker. It reads the
detector scores and the settings that the tracker declares, and it decides
while running whether the tracker needs help. The main idea is simple: the
detector's own scores show whether its output fits the tracker. \scl{} looks
at the scores of recent frames and checks how clearly they separate confident
objects from clutter. If the separation is clear, the tracker keeps its own
settings, and the layer only adds help that the tracker lacks, such as a step
that continues tracks. If it is not clear, \scl{} chooses which candidates
reach the tracker, assigns them new scores, and sets the tracker's thresholds.
The novelty is therefore a layer outside the tracker that needs no labels,
uses one fixed configuration for many trackers, and stays out of the way when
the tracker already fits the detector.

The main contributions are:
\begin{enumerate}
\item \scl{}, a control layer that needs no training, adjusts itself from the
      detector scores, uses no future frames, keeps the tracker's settings
      when the scores are clean, and takes full control only when they are
      noisy (Sect.~\ref{sec:method}).
\item A host contract, a set of four numbers from the tracker's published
      settings, that lets one fixed configuration work with single-stage and
      two-stage trackers without changing them (Sect.~\ref{sec:contract}).
\item A development path from always-on banding (\aob{}), an earlier design
      that always intervened, in which each step is compared with that
      baseline (Sect.~\ref{sec:development}).
\item A test of one fixed configuration on seven published trackers, four
      detection sources, and four datasets, with development and predeclared
      external groups, paired bootstrap intervals, and a speed measurement on
      a central processing unit (CPU) (Sects.~\ref{sec:setup} and~\ref{sec:results}).
\end{enumerate}

The rest of the paper is organized as follows. Section~\ref{sec:data-eval}
describes the data, the evaluation roles, and the metrics, and
Section~\ref{sec:related} reviews related work. Section~\ref{sec:method}
explains \scl{} and all its settings, and Section~\ref{sec:development} shows
how it was developed from an earlier design that always intervened. The
configuration was then frozen before any external system was run.
Section~\ref{sec:setup} gives the experimental setup,
Section~\ref{sec:results} the results and their analysis, and
Section~\ref{sec:conclusion} concludes the paper.
